\section*{Conclusions}\label{chap:conclusion}
\addcontentsline{toc}{section}{Conclusions}

In this analysis, a search for a narrow \(Z^{\prime} \to t\bar{t}\) resonance was performed using \(1\,\mathrm{fb}^{-1}\) of \(\sqrt{s} = 13\,\mathrm{TeV}\) data from the ATLAS detector. The reconstructed \(t\bar{t}\) invariant mass \(m_{t\bar{t}}\) was selected as the final discriminant due to its superior separation power between signal and background. After applying optimized event selection criteria, a binned \(\chi^2\) test of the \(m_{t\bar{t}}\) distribution was applied to evaluate the agreement between background and the observed data. That resulted in a \(p\)-value of $1.34 \times 10^{-2}$ and thus does not provide
evidence for a signal. However, the contribution of systematic uncertainties was not taken into account. After including their contribution, a \(p\)-value of 0.998 was obtained, showing no significant deviation from Standard Model expectations.

As no significant excess above the Standard Model background prediction was observed, exclusion limits at the 95\% confidence level were derived on the $Z'$ production cross section multiplied by its branching fraction to top-quark pairs, $\sigma(pp \to Z') \times \mathcal{B}(Z' \to t\bar{t})$. These limits were presented as a function of the hypothesized $Z'$ mass, showing both the theoretical prediction and the 95\% CL exclusion limits on the production cross-section. It was found that, for $Z'$ masses
up to several TeV, production rates above a few picobarns are excluded.

